% Part: intuitionistic-logic
% Chapter: tableaux
% Section: rules

\documentclass[../../../include/open-logic-section]{subfiles}

\begin{document}

\olfileid{int}{tab}{rul}

\olsection{સ્ફુરણાવાદી તર્ક માટેના નિયમો}

$\land$ અને $\lor$ માટેના નિયમો સામાન્ય વિધાનસંબંધી
ચિહ્નિત !!{tableau}s જેવા જ છે; ફક્ત ઉપસર્ગો ઉમેરાય છે.
દરેક કિસ્સામાં ચિહ્નિત !!{formula}
$\sFmla{S}{!A}[\sigma]$ પર લાગુ થયેલો નિયમ નવાં
!!{formula}s આપે છે, જેમને પણ~$\sigma$ ઉપસર્ગ હોય છે.
આ સહજ છે: દાખલા તરીકે, $!A \land !B$ એ~$\sigma$
નામના જગતમાં સત્ય હોય, તો $!A$ અને $!B$ પણ
~$\sigma$માં સત્ય છે. આ નિયમના નિષ્કર્ષ પર અન્ય
જગતનો ઉપસર્ગ મૂકાતો નથી.
\sourcecorrection{OLMD-098}{મૂળનું કૌંસ કહે છે કે
$!A$ અને $!B$ બીજા કોઈ જગતમાં સત્ય નથી;
સ્ફુરણાવાદી એકદિશવર્ધિતાથી તેઓ સુલભ જગતોમાં પણ
સત્ય હોય છે. નિયમ માત્ર એ જ ઉપસર્ગ પર નિષ્કર્ષ
લખે છે એમ સ્પષ્ટ કર્યું છે.}
$\land$ અને $\lor$ના નિયમો \olref{tab:prop-rules}માં
એકત્ર કર્યા છે.

\begin{table}
  \[\def\arraystretch{3}\begin{array}{|c|c|}
    \hline
    \AxiomC{\sFmla{\True}{!A \land !B}[\sigma]}
    \RightLabel{\TRule{\True}{\land}}
    \UnaryInfC{\sFmla{\True}{!A}[\sigma]}
    \noLine
    \UnaryInfC{\sFmla{\True}{!B}[\sigma]}
    \DisplayProof
    &
    \AxiomC{\sFmla{\False}{!A \land !B}[\sigma]}
    \RightLabel{\TRule{\False}{\land}}
    \UnaryInfC{$\sFmla{\False}{!A}[\sigma] \quad \mid \quad
      \sFmla{\False}{!B}[\sigma]$}
    \DisplayProof
    \\[2ex]
    \hline
    \AxiomC{\sFmla{\True}{!A \lor !B}[\sigma]}
    \RightLabel{\TRule{\True}{\lor}}
    \UnaryInfC{$\sFmla{\True}{!A}[\sigma] \quad \mid \quad
      \sFmla{\True}{!B}[\sigma]$}
    \DisplayProof
    &
    \AxiomC{\sFmla{\False}{!A \lor !B}[\sigma]}
    \RightLabel{\TRule{\False}{\lor}}
    \UnaryInfC{\sFmla{\False}{!A}[\sigma]}
    \noLine
    \UnaryInfC{\sFmla{\False}{!B}[\sigma]}
    \DisplayProof
    \\[2ex]
    \hline
  \end{array}\]
  \caption{$\land$ અને $\lor$ માટે ઉપસર્ગવાળા !!{tableau} નિયમો}
  \ollabel{tab:prop-rules}
\end{table}

બંધતાની શરત સામાન્ય !!{tableau}s જેવી છે, પરંતુ
!!{formula}s સાથે ઉપસર્ગો પણ સુસંગત હોવા જોઈએ. વધુમાં,
અહીં થોડી વ્યાપક શરત લઈ શકીએ: સ્ફુરણાવાદી
નિદર્શોમાં !!{formula}s એક વાર સત્ય બને, પછી
સત્ય રહે છે. તેથી $!A$ એ~$\sigma$માં સત્ય હોય
અને સુલભ ઉપસર્ગ~$\sigma.{*}$ પર ખોટું હોય
એવું શક્ય નથી. આથી કોઈ શાખામાં નીચેના બંને હોય,
તો તે બંધ છે:
\[
\sFmla{\True}{!A}[\sigma] \quad\text{અને}\quad \sFmla{\False}{!A}[\sigma.{*}]
\]
કોઈ ઉપસર્ગ $\sigma$ અને !!{formula}~$!A$ માટે.
નોંધો કે ચિહ્નો ઊલટાં હોય, એટલે શાખામાં
\[
\sFmla{\False}{!A}[\sigma] \quad\text{અને}\quad \sFmla{\True}{!A}[\sigma.{*}]
\]
હોય, તો $*$ ખાલી શ્રેણી હોય ત્યારે જ શાખા બંધ છે.

આ ઉપરાંત, શાખામાં~$\sFmla{\True}{\bot}[\sigma]$
હોય, તો તે બંધ છે.

ધારણાઓ ગોઠવવાના નિયમો પણ સામાન્ય !!{tableau}s
જેવા જ છે, પરંતુ ધારણાઓ માટે આપણે હંમેશા
ઉપસર્ગ~$1$ લઈએ છીએ. (કયો ઉપસર્ગ લઈએ તે
મહત્ત્વનું નથી, જો બધી ધારણાઓ માટે તે એક જ હોય.)
દાખલા તરીકે, આપણે કહીએ છીએ કે
\[
!B_1, \dots, !B_n \Proves !A
\]
ત્યારે અને માત્ર ત્યારે જ જ્યારે નીચેની ધારણાઓ
માટે બંધ ટેબ્લો હોય:
\[
\sFmla{\True}{!B_1}[1], \dots, \sFmla{\True}{!B_n}[1],
\sFmla{\False}{!A}[1].
\]

શરતી~$\lif$ માટેના નિયમો પ્રશિષ્ટ અને મોડલ
કિસ્સાઓથી જુદા છે. $\TRule{\True}{\lif}$ નિયમ
$\sFmla{\True}{!A \lif !B}[\sigma]$ ધરાવતી શાખાને
જુદી બે શાખાઓ પર અનુક્રમે
$\sFmla{\False}{!A}[\sigma.{*}]$ અને
$\sFmla{\True}{!B}[\sigma.{*}]$ ઉમેરીને વિસ્તારે છે.
\sourcecorrection{OLMD-099}{મૂળ ગદ્યમાં આ નિયમનાં
બંને નિષ્કર્ષોના સત્યચિહ્ન ઊલટાં છે: ત્યાં
$\sFmla{\True}{!A}[\sigma.{*}]$ અને
$\sFmla{\False}{!B}[\sigma.{*}]$ છે, જ્યારે
નીચેની નિયામક કોષ્ટકમાં સાચાં
$\sFmla{\False}{!A}[\sigma.{*}]$ અને
$\sFmla{\True}{!B}[\sigma.{*}]$ છે. મૂળના
બંને $\TRule$ ગદ્ય-લેબલોમાં દલીલોનો ક્રમ પણ
કોષ્ટકથી ઊલટો છે; અહીં કોષ્ટકનો ક્રમ લીધો છે.}
આ નિયમ ફક્ત એવા ઉપસર્ગ~$\sigma.{*}$ માટે
લાગુ કરી શકાય જે સંબંધિત શાખા પર \emph{પહેલેથી}
આવે છે. આવા ઉપસર્ગને તે શાખા પર ``વપરાયેલો''
કહીએ. ($\sigma.{*}$માં $\sigma$ પોતે પણ આવે છે,
તેથી જુદી શાખાઓ પર અનુક્રમે
$\sFmla{\False}{!A}[\sigma]$ અને
$\sFmla{\True}{!B}[\sigma]$ ઉમેરીને આ નિયમ
હંમેશા લાગુ કરી શકાય છે.)

$\TRule{\False}{\lif}$ નિયમ
$\sFmla{\False}{!A \lif !B}[\sigma]$ ધરાવતી શાખાને
એ જ શાખા પર $\sFmla{\True}{!A}[\sigma.n]$
અને $\sFmla{\False}{!B}[\sigma.n]$ બંને ઉમેરીને
વિસ્તારે છે; અહીં $\sigma.n$ એ શાખા પર નવો
ઉપસર્ગ છે.

$\lnot$ માટેના નિયમો સમાન રીતે વ્યાખ્યાયિત થાય છે
($\lnot !A$ને $!A \lif \lfalse$ તરીકે વ્યાખ્યાયિત
કર્યા પ્રમાણે).

નિયમો \olref{tab:rules-lif-lnot}માં આપ્યા છે.

\begin{table}
  \[\def\arraystretch{3}\begin{array}{|c|c|}
    \hline
    \AxiomC{\sFmla{\True}{\lnot !A}[\sigma]}
    \RightLabel{\TRule{\True}{\lnot}}
    \UnaryInfC{$\sFmla{\False}{!A}[\sigma.{*}]$}
    \DisplayProof
    &
    \AxiomC{\sFmla{\False}{\lnot !A}[\sigma]}
    \RightLabel{\TRule{\False}{\lnot}}
    \UnaryInfC{\sFmla{\True}{!A}[\sigma.n]}
    \DisplayProof
    \\[1ex]
    \text{$\sigma.{*}$ વપરાયેલો છે} & \text{$\sigma.n$ નવો છે}\\
    \hline
    \AxiomC{\sFmla{\True}{!A \lif !B}[\sigma]}
    \RightLabel{\TRule{\True}{\lif}}
    \UnaryInfC{$\sFmla{\False}{!A}[\sigma.{*}] \quad \mid \quad
      \sFmla{\True}{!B}[\sigma.{*}]$}
    \DisplayProof
    &
    \AxiomC{\sFmla{\False}{!A \lif !B}[\sigma]}
    \RightLabel{\TRule{\False}{\lif}}
    \UnaryInfC{\sFmla{\True}{!A}[\sigma.n]}
    \noLine
    \UnaryInfC{\sFmla{\False}{!B}[\sigma.n]}
    \DisplayProof
    \\[1ex]
    \text{$\sigma.{*}$ વપરાયેલો છે} & \text{$\sigma.n$ નવો છે}\\
    \hline
  \end{array}\]
  \caption{$\lnot$ અને $\lif$ માટે ઉપસર્ગવાળા !!{tableau} નિયમો}
  \ollabel{tab:rules-lif-lnot}
\end{table}

\end{document}
